% Level 3, Week 5 answers -- Sorting, Stacks and Queues
\answersweek{5}{Sorting, Stacks and Queues}{Answers for workbook chapter 5 and homework 5.}

\section{Workbook exercises}
\subsection{Exercise 5.1: Sort Race}
\pyfile{../excercises/week05/ex-sort-race/solution.py}
\begin{answer}[Exercise 5.1: Sort Race]
\textbf{Swaps for the 8 cards} \texttt{[7, 3, 9, 1, 5, 3, 8, 2]}: selection 5, bubble 16. Both give \texttt{[1, 2, 3, 3, 5, 7, 8, 9]} and the original list is unchanged.\\
\textbf{What to look for:} the tuple swap; \texttt{swaps} counted only when a swap happens; a copy made with \texttt{list(items)} before sorting.
\end{answer}

\subsection{Exercise 5.2: Leaderboard}
\pyfile{../excercises/week05/ex-leaderboard/solution.py}
\begin{answer}[Exercise 5.2: Leaderboard]
Expected board: \texttt{1. Mei - 150}, \texttt{2. Ada - 120}, \texttt{3. Leo - 95}, \texttt{3. Omar - 95}, \texttt{5. Zoe - 80}.\\
\textbf{What to look for:} \texttt{key=points\_of} without brackets; \texttt{[:n]} to keep the top \texttt{n}; the tie rule changes \texttt{place} only when the points differ from the previous entry. A \texttt{lambda} key is fine for students who know it.
\end{answer}

\subsection{Exercise 5.3: Undo Paint}
\pyfile{../excercises/week05/ex-undo-paint/solution.py}
\begin{answer}[Exercise 5.3: Undo Paint]
\textbf{What to look for:} \texttt{push} appends to the end; \texttt{pop} and \texttt{peek} return \texttt{None} on an empty stack instead of crashing; \texttt{peek} does not remove anything.
In the program, pressing \texttt{u} repeatedly removes strokes newest first, and pressing it on an empty canvas does nothing.
The colour name is written at the top, so players do not have to rely on colour alone.
\end{answer}

\section{Homework}
\subsection{Core: Turn Queue}
\pyfile{../excercises/week05/hw-turn-queue/solution.py}
\begin{answer}[Core: Turn Queue]
\textbf{Prediction:} Priya, Sam, Kai, Zoe, Priya, Sam. After Priya plays and Kai leaves, the queue is Sam, Zoe, Priya.
\end{answer}

\begin{answer}[Challenge: Stack or queue?]
Back button: \textbf{S} (the last page you visited is the first you go back to). Playlist: \textbf{Q} (songs play in the order they were added). Printing: \textbf{Q} (first sent, first printed). Pile of homework, newest on top: \textbf{S} (you take the top, newest sheet first).
Accept other answers with a sensible reason (for example a \enquote{shuffle} playlist is neither).
\end{answer}

\begin{answer}[Creative: Design a fair leaderboard]
\textbf{Good looks like:} nicknames only; ties share a place; at least one effort-based award; ideally a way to hide yourself from the board. Ask the student to explain who their design is kind to.
\end{answer}
